\usepackage{array}
\usepackage{shadowtext}
\shadowoffset{1pt}
\shadowcolor{black}

\newcolumntype{L}[1]{>{\raggedright\arraybackslash}p{#1}}
\newcolumntype{R}[1]{>{\raggedleft\arraybackslash}p{#1}}
\newcolumntype{C}[1]{>{\centering\arraybackslash}p{#1}}




\newcommand{\cuttingBoardBackground}[9]{
 \begin{tikzpicture}[x=\linewidth,y=0.75\linewidth] % Set aspect ratio to 4:3
    % Draw the cutting board background
    \node (image) [anchor=south west,inner sep=0] (image) at (0,0) {
        \includegraphics[width=\linewidth,trim={15.5cm 0 15.5cm 0}, clip, keepaspectratio]{img/snijplank.jpg}
    };
    \begin{scope}[x={(image.south east)},y={(image.north west)}]
        \def\h{2.5}
        \node (img1) at (image.north west) [xshift=2cm, yshift=-2cm] {\includegraphics[height=\h cm]{img/cards/generated/#1.png}};
        \node (img2) at (img1.east) [xshift=\h*0.25cm] {\includegraphics[height=\h cm]{img/cards/generated/#2.png}};
    

   % Add the new title
   \node[anchor=north west, inner sep=0] at (img1.east) [yshift=-2cm, xshift=0cm] {\shadowtext{\textbf{Nieuwe stapels}}
   };
   


    \end{scope}

 \end{tikzpicture}
}

\renewcommand{\thefootnote}{\arabic{footnote}}
\newcommand{\footnoteFritsjeDesNemen}{\textbf{een Fritsje des nemen} is equivalent aan \Fritsen (zie definities \ref{definitie:fritsen_equivalent} en \ref{definitie:fritsen} op pagina \pageref{definitie:fritsen})}
\newcommand{\footnoteEenDubbeleFritsNemen}{\textbf{een dubbele Frits nemen} is equivalent aan twee keer \Fritsen (zie definities \ref{definitie:fritsen} en \ref{definitie:dubbelfritsen} op pagina \pageref{definitie:fritsen})}
\newcommand{\footnoteDubbelFritsVariantToevoeging}{\textbf{Dubbelfritsen} (zie definitie \ref{definitie:dubbelfritsen} op pagina \pageref{definitie:dubbelfritsen})}

\newlength\q
\setlength\q{\dimexpr .363\textwidth -4\tabcolsep}

\newcommand{\drawBar}[1]{
\begin{minipage}[c][0.5cm][c]{1.0\textwidth}
    \hspace*{-0.75cm} 
    \ifodd\thepage\else\hfill\fi
    \textbf{\thepage} / \textbf{\reallastpage}
    \ifodd\thepage\else\hspace*{0.55cm}\fi
\end{minipage}
}

\newcommand{\drawParagraph}[1]{
\paragraph{#1} \texttt{ } \newline
}

\newcommand{\createlabel}[2]{%
    \expandafter\edef\csname#1\endcsname{#2}
}

\newcommand{\EenSpeler}{\textbf{[Een speler] }}
\newcommand{\EenSpelerN}{\textbf{[Een speler]}}

\newcommand{\eenSpeler}{\textbf{[een speler] }}
\newcommand{\eenSpelerN}{\textbf{[een speler]}}

\newcommand{\huidigeSpeler}{\textbf{[huidige speler] }}
\newcommand{\huidigeSpelerN}{\textbf{[huidige speler]}}

\newcommand{\HuidigeSpeler}{\textbf{[Huidige speler] }}
\newcommand{\HuidigeSpelerN}{\textbf{[Huidige speler]}}

\newcommand{\andereSpeler}{\textbf{[andere speler] }}
\newcommand{\andereSpelerN}{\textbf{[andere speler]}}

\newcommand{\AndereSpeler}{\textbf{[Andere speler] }}
\newcommand{\AndereSpelerN}{\textbf{[Andere speler]}}

\newcommand{\andereSpelers}{\textbf{[andere spelers] }}
\newcommand{\andereSpelersN}{\textbf{[andere spelers]}}

\newcommand{\AndereSpelers}{\textbf{[Andere spelers] }}
\newcommand{\AndereSpelersN}{\textbf{[Andere spelers]}}

\newcommand{\alleSpelers}{\textbf{[alle spelers]} }
\newcommand{\alleSpelersN}{\textbf{[alle spelers]}}

\newcommand{\AlleSpelers}{\textbf{[Alle spelers] }}
\newcommand{\AlleSpelersN}{\textbf{[Alle spelers]}}

\newcommand{\MedeSpeler}{\textbf{[Medespeler] }}
\newcommand{\MedeSpelerN}{\textbf{[Medespeler]}}

\newcommand{\medeSpeler}{\textbf{[medespeler]} }
\newcommand{\medeSpelerN}{\textbf{[medespeler]}}

\newcommand{\MedeSpelers}{\textbf{[Medespelers] }}
\newcommand{\MedeSpelersN}{\textbf{[Medespelers]}}

\newcommand{\medeSpelers}{\textbf{[medespelers]} }
\newcommand{\medeSpelersN}{\textbf{[medespelers]}}

\newcommand{\VorigeSpeler}{\textbf{[Vorige speler] }}
\newcommand{\VorigeSpelerN}{\textbf{[Vorige speler]}}

\newcommand{\vorigeSpeler}{\textbf{[vorige speler] }}
\newcommand{\vorigeSpelerN}{\textbf{[vorige speler]}}

\newcommand{\VolgendeSpeler}{\textbf{[Volgende speler] }}
\newcommand{\VolgendeSpelerN}{\textbf{[Volgende speler]}}

\newcommand{\volgendeSpeler}{\textbf{[volgende speler] }}
\newcommand{\volgendeSpelerN}{\textbf{[volgende speler]}}

\newcommand{\deSpeler}{\textbf{[de speler] }}
\newcommand{\deSpelerN}{\textbf{[de speler]}}

\newcommand{\DeSpeler}{\textbf{[De speler] }}
\newcommand{\DeSpelerN}{\textbf{[De speler]}}

\newcommand{\Willem}{\textbf{[Willem] }}
\newcommand{\WillemN}{\textbf{[Willem]}}

\newcommand{\Frits}{\textbf{[Frits] }}
\newcommand{\FritsN}{\textbf{[Frits]}}

\newcommand{\Lisa}{\textbf{[Lisa] }}
\newcommand{\LisaN}{\textbf{[Lisa]}}

\newcommand{\Kim}{\textbf{[Kim] }}
\newcommand{\KimN}{\textbf{[Kim]}}

\newcommand{\Fritsen}{\textbf{Fritsen} }
\newcommand{\FritsenN}{\textbf{Fritsen}}

\newcommand{\customBox}[1]{
    \begin{framed}
    \noindent
    #1
    \end{framed}
}

\newcommand{\customBoxItalic}[1]{
    \begin{framed}
    \noindent
    \textit{#1}
    \end{framed}
}


\newcommand{\kaart}[1]{\textbf{\ul{#1}}}

\newcommand{\hoofdstuk}[2]{\vspace{-0.5cm}\section*{#1. #2}}
\newcommand{\deelhoofdstuk}[1]{\vspace{-0.5cm}\noindent\subsection*{#1}}

\newcommand{\beginLijst}[1]{
    \def \sectieNummer {#1}
    \begin{enumerate}[label=\texttt{\sectieNummer-\arabic*},align=left, labelwidth=31pt, leftmargin=37pt]
}

\newcommand{\vervolgLijst}[0]
    {\begin{enumerate}[label=\texttt{\sectieNummer-\arabic*},resume,align=left, labelwidth=31pt, leftmargin=37pt]
}

\newcommand{\beginLijstKlein}[1]{
    \def \sectieNummer {#1}
    \begin{enumerate}[label=\texttt{\sectieNummer-\arabic*},align=left, labelwidth=21pt, leftmargin=27pt]
}

\newcommand{\vervolgLijstKlein}[0]
    {\begin{enumerate}[label=\texttt{\sectieNummer-\arabic*},resume,align=left, labelwidth=21pt, leftmargin=27pt]
}

\newcommand{\numeriekeLijst}[0]
    {\begin{enumerate}[label=\arabic*.,nolistsep]
}

\newcommand{\eindNumeriekeLijst}[0]{
    \end{enumerate}
}

\newcommand{\beginABCLijst}[1]{
    \def \sectieNummer {#1}
    \begin{enumerate}[label=\texttt{\sectieNummer\Alph*},nolistsep,align=left, labelwidth=25pt, leftmargin=31pt]
}

\newcommand{\beginABCLijstKlein}[1]{
    \def \sectieNummer {#1}
    \begin{enumerate}[label=\texttt{\sectieNummer\Alph*},nolistsep,align=left, labelwidth=13pt, leftmargin=19pt]
}


\newcommand{\eindABCLijst}[0]{
    \end{enumerate}
}

\newcommand{\puntLijst}[0]{
    \begin{itemize}[label=$\circ$,nolistsep]
}

\newcommand{\eindPuntLijst}[0]{
    \end{itemize}
}

\newcommand{\eindLijst}[0]{
    \end{enumerate}
}

\newcommand{\quotes}[1]{\textquotedblleft\textit{#1}\textquotedblright}

\setlength\tabcolsep{0pt}
\newcommand{\kaartnaam}[5]{
\begin{minipage}[c][2.25cm][c]{1.0\textwidth}
    \begin{tabularx}{\textwidth}{Xccccm{0.00cm}}
        \vspace{-1.6cm}
        \begin{flushleft}
            #1
        \end{flushleft}&
        \ifx&#2&%
            &
        \else
            \includegraphics[height=2cm]{img/cards/generated/#2.png} \hspace{0.05cm} &
        \fi
        \ifx&#3&%
            &
        \else
            \includegraphics[height=2.0cm]{img/cards/generated/#3.png} \hspace{0.05cm} &
        \fi
        \ifx&#4&%
            &
        \else
            \includegraphics[height=2.0cm]{img/cards/generated/#4.png} \hspace{0.05cm} &
        \fi
        \ifx&#5&%
            &
        \else
            \includegraphics[height=2.0cm]{img/cards/generated/#5.png} &
        \fi
    \end{tabularx}
\end{minipage}
}

\newcommand{\zetKort}[8]{
\deelhoofdstuk{#1}
    \noindent
    \begin{tikzpicture}
        \def\h{2.5}
        \node (img1) {\includegraphics[height=\h cm]{img/cards/generated/#2.png}};
        \node (img2) at (img1.south east) [xshift=-\h*0.29cm,yshift=\h*0.455cm] {\includegraphics[height=\h cm]{img/cards/generated/#3.png}};
    \end{tikzpicture}
    \begin{tikzpicture}
        \def\h{2.5}
        \node (img1) {\includegraphics[height=\h cm]{img/cards/generated/#4.png}};
        \node (img2) at (img1.south east) [xshift=-\h*0.29cm,yshift=\h*0.455cm] {\includegraphics[height=\h cm]{img/cards/generated/#5.png}};
    \end{tikzpicture}
    \begin{tikzpicture}
        \def\h{2.5}
        \node (img1) {\includegraphics[height=\h cm]{img/cards/generated/#6.png}};
        \node (img2) at (img1.south east) [xshift=-\h*0.29cm,yshift=\h*0.455cm] {\includegraphics[height=\h cm]{img/cards/generated/#7.png}};
    \end{tikzpicture}
    \newline
    #8
}

\newcommand{\zetKortMetKruis}[8]{
\deelhoofdstuk{#1}
    \noindent
    \begin{tikzpicture}
        \def\h{2.5}
        \node (img1) {\includegraphics[height=\h cm]{img/cards/generated/#2.png}};
        \node (img2) at (img1.south east) [xshift=-\h*0.29cm,yshift=\h*0.455cm] {\includegraphics[height=\h cm]{img/cards/generated/#3.png}};
    \end{tikzpicture}
    \begin{tikzpicture}
        \def\h{2.5}
        \node (img1) {\includegraphics[height=\h cm]{img/cards/generated/#4.png}};
        \node (img2) at (img1.south east) [xshift=-\h*0.29cm,yshift=\h*0.455cm] {\includegraphics[height=\h cm]{img/cards/generated/#5.png}};
    \end{tikzpicture}
    \begin{tikzpicture}
        \def\h{2.5}
        \node (img1) {\includegraphics[height=\h cm]{img/cards/generated/#6.png}};
        \node (img2) at (img1.south east) [xshift=-\h*0.29cm,yshift=\h*0.455cm] {\includegraphics[height=\h cm]{img/cards/generated/#7.png}};
        
        \draw [red,ultra thick,rounded corners]([xshift=\h*0.04cm,yshift=-\h*0.04cm]img1.north west) -- ([xshift=-\h*0.04cm,yshift=\h*0.04cm]img2.south east);
        \draw [red,ultra thick,rounded corners]([xshift=-\h*0.04cm,yshift=\h*0.05cm]img2.north east) -- ([xshift=-\h*0.085cm,yshift=\h*0.04cm]img2.south west);   
    \end{tikzpicture}
    \newline
    #8
}

\newcommand{\zetKoningAas}[5]{
    \noindent
    \begin{minipage}[t][3.5cm][t]{1.0\textwidth}
        \begin{wrapfigure}{L}{10.9cm}
            \vspace{-15pt}
            \begin{tikzpicture}
                \def\h{2.4}
                \node (img1) {\includegraphics[height=\h cm]{img/cards/generated/queen_of_#1.png}};
                \node (img2) at (img1.south east) [xshift=-\h*0.27cm,yshift=\h*0.455cm] {\includegraphics[height=\h cm]{img/cards/generated/king_of_#1.png}};
                \node (img3) at (img2.south east) [xshift=-\h*0.27cm,yshift=\h*0.555cm] {\includegraphics[height=\h cm]{img/cards/generated/ace_of_#1.png}};
            \end{tikzpicture}
            \begin{tikzpicture}
                \def\h{2.4}
                \node (img1) {\includegraphics[height=\h cm]{img/cards/generated/queen_of_#2.png}};
                \node (img2) at (img1.south east) [xshift=-\h*0.27cm,yshift=\h*0.455cm] {\includegraphics[height=\h cm]{img/cards/generated/king_of_#2.png}};
                \node (img3) at (img2.south east) [xshift=-\h*0.27cm,yshift=\h*0.555cm] {\includegraphics[height=\h cm]{img/cards/generated/ace_of_#2.png}};
            \end{tikzpicture}
            \begin{tikzpicture}
                \def\h{2.4}
                \node (img1) {\includegraphics[height=\h cm]{img/cards/generated/queen_of_#3.png}};
                \node (img2) at (img1.south east) [xshift=-\h*0.27cm,yshift=\h*0.455cm] {\includegraphics[height=\h cm]{img/cards/generated/king_of_#3.png}};
                \node (img3) at (img2.south east) [xshift=-\h*0.27cm,yshift=\h*0.555cm] {\includegraphics[height=\h cm]{img/cards/generated/ace_of_#32.png}};
            \end{tikzpicture}
            \begin{tikzpicture}
                \def\h{2.4}
                \node (img1) {\includegraphics[height=\h cm]{img/cards/generated/queen_of_#4.png}};
                \node (img2) at (img1.south east) [xshift=-\h*0.27cm,yshift=\h*0.455cm] {\includegraphics[height=\h cm]{img/cards/generated/king_of_#4.png}};
                \node (img3) at (img2.south east) [xshift=-\h*0.27cm,yshift=\h*0.555cm] {\includegraphics[height=\h cm]{img/cards/generated/ace_of_#4.png}};
            \end{tikzpicture}
        \end{wrapfigure}
        #5.
    \end{minipage}
}

\newcommand{\zetLangMetKruis}[9]{
    \noindent
    \begin{minipage}[t][#1cm][t]{1.0\textwidth}
        \begin{wrapfigure}{L}{5.0cm}
            \vspace{-15pt}
            \begin{tikzpicture}
                \def\h{2.5}
                \node (img1) {\includegraphics[height=\h cm]{img/cards/generated/#2.png}};
                \node (img2) at (img1.south east) [xshift=-\h*0.27cm,yshift=\h*0.455cm] {\includegraphics[height=\h cm]{img/cards/generated/#3.png}};
            \end{tikzpicture}
            \begin{tikzpicture}
                \def\h{2.5}
                \node (img1) {\includegraphics[height=\h cm]{img/cards/generated/#4.png}};
                \node (img2) at (img1.south east) [xshift=-\h*0.27cm,yshift=\h*0.455cm] {\includegraphics[height=\h cm]{img/cards/generated/#5.png}};
                
                \draw [red,ultra thick,rounded corners]([xshift=\h*0.04cm,yshift=-\h*0.04cm]img1.north west) -- ([xshift=-\h*0.04cm,yshift=\h*0.04cm]img2.south east);
                \draw [red,ultra thick,rounded corners]([xshift=-\h*0.04cm,yshift=\h*0.05cm]img2.north east) -- ([xshift=-\h*0.085cm,yshift=\h*0.04cm]img2.south west);   
            \end{tikzpicture}
        \end{wrapfigure}
        #6: 
        \ifx&#7&%
        \else
            \ifx&#8&%
                \begin{itemize}[nolistsep,leftmargin=*]
                    \item #7
                \end{itemize}       
            \else
                \begin{itemize}[nolistsep,leftmargin=*]
                    \item #7
                    \item #8
                \end{itemize}
            \fi
        \fi
        \ifx&#9&%
        \else
            \ifx&#7&%
                \newline
                #9
            \else
                #9
            \fi
        \fi
    \end{minipage}
}

\newcommand{\zetLang}[9]{
    \noindent
    \begin{minipage}[t][#1cm][t]{1.0\textwidth}
        \begin{wrapfigure}{L}{5.0cm}
            \vspace{-15pt}
            \begin{tikzpicture}
                \def\h{2.5}
                \node (img1) {\includegraphics[height=\h cm]{img/cards/generated/#2.png}};
                \node (img2) at (img1.south east) [xshift=-\h*0.27cm,yshift=\h*0.455cm] {\includegraphics[height=\h cm]{img/cards/generated/#3.png}};
            \end{tikzpicture}
            \begin{tikzpicture}
                \def\h{2.5}
                \node (img1) {\includegraphics[height=\h cm]{img/cards/generated/#4.png}};
                \node (img2) at (img1.south east) [xshift=-\h*0.27cm,yshift=\h*0.455cm] {\includegraphics[height=\h cm]{img/cards/generated/#5.png}}; 
            \end{tikzpicture}
        \end{wrapfigure}
        #6: 
        \ifx&#7&%
        \else
            \ifx&#8&%
                \begin{itemize}[nolistsep,leftmargin=*]
                    \item #7
                \end{itemize}       
            \else
                \begin{itemize}[nolistsep,leftmargin=*]
                    \item #7
                    \item #8
                \end{itemize}
            \fi
        \fi
        \ifx&#9&%
        \else
            \ifx&#7&%
                \newline
                #9
            \else
                #9
            \fi
        \fi
    \end{minipage}
}

\newcommand{\proLabel}[0]{
    \hspace{-0.25cm}\tcbox{\textit{\textcolor{white}{PRO}}}
}

\newcommand{\proLabelMetRuimte}[0]{
    \tcbox{\textit{\textcolor{white}{PRO}}}
}

\newcommand{\tweedeDruk}[0]{
\begingroup
    \fontsize{30pt}{36pt}\selectfont
    \tcbox{\textit{\textcolor{white}{2\hspace*{0.1cm}\textsuperscript{e} druk}}}
\endgroup
}




\def\Put(#1,#2)#3{\leavevmode\makebox(0,0){\put(#1,#2){#3}}}
